\section{Proofs for the complete-fleet statements}\label{app:theory-proofs}

\begin{proof}[Proof of Proposition~\ref{prop:theory-support}]
For physical $x$, $p\in\partial F(e(x))$, and any hull point $(e,c)$,
linear minimization over the convex hull gives
$c+p^\top e\geq V(p)$.  The subgradient inequality gives
$F(e)\geq F(e(x))+p^\top(e-e(x))$.  Adding and minimizing over the hull
yields $CH\geq c(x)+F(e(x))-r(x;p)$, proving the first inequality;
$c(x)+F(e(x))\geq D\geq CH$ proves the second.  If $r(x;p)=0$, all
inequalities collapse to $CH=D=c(x)+F(e(x))$.  Conversely, compactness
supplies a physical minimizer when $D=CH$.  Under the relative-interior
condition, Fenchel--Rockafellar duality gives an attained maximum
$CH=V(p^*)-F^*(p^*)$.  Under the alternative minorant condition, apply
the same duality theorem to that finite minorant.  It agrees with $F$
on all hull loads, while its conjugate bounds $F^*$ from above; weak
duality for $F$ then makes the same $p^*$ optimal for $F$.  For any
physical minimizer $x^*$, the two
nonnegative Fenchel terms satisfy
\[
r(x^*;p^*)+
F(e(x^*))-(p^*)^\top e(x^*)+F^*(p^*)
=D-CH=0.
\]
Each term is therefore zero.  Fenchel equality gives
$p^*\in\partial F(e(x^*))$ and the first equality gives
$r(x^*;p^*)=0$.  This holds for every physical minimizer and every
attained hull-dual maximizer.
\end{proof}

\begin{proof}[Proof of Corollary~\ref{cor:best-price}]
Let $(e^*,c^*)$ be a hull minimizer with $e^*=e(x)$, and let $p^*$ be an
attained hull-dual maximizer. The nonnegative terms
$c^*+(p^*)^\top e^*-V(p^*)$ and
$F(e^*)-(p^*)^\top e^*+F^*(p^*)$ sum to zero by strong duality.
Hence $p^*\in\partial F(e^*)$ and
$r(x;p^*)=c(x)-c^*=c(x)+F(e(x))-CH=H(x)$.
Proposition~\ref{prop:theory-support} supplies the matching lower bound
for every price in $\partial F(e(x))$.
\end{proof}

For the accounting identity, $F^*(p)\geq p^\top e(x)-F(e(x))$ by
definition, so supplier LOC is nonnegative.  Substitution in
\eqref{eq:theory-loc} gives the claimed sum.  The same
qualification above yields an attained hull-dual price
$p^*$ and $V(p^*)-F^*(p^*)=CH$.  At a physical planner optimizer the
LOC sum is $D-CH$.  At any price in the plan's subdifferential, Fenchel
equality makes supplier LOC zero.

\begin{proof}[Proof of \eqref{eq:theory-upper}]
Let $(e^*,c^*)$ be a hull minimizer, $p^*=a+Qe^*$,
$d=\|e(x)-e^*\|_Q$, and $H=H(x)$.  First-order hull optimality makes
$(e^*,c^*)$ a minimizer of $c+(p^*)^\top e$, so
\[
 \mathrm{LOC}_{\rm fleet}(x;p^*)=H-\tfrac12d^2\geq0.
\]
In particular $d\leq\sqrt{2H}$.  For each physical alternative $y$,
the change in its linear advantage when the price moves from $p^*$ to
$p_x=p^*+Q(e(x)-e^*)$ is
$(e(x)-e^*)^\top Q(e(x)-e(y))$, at most $dR_Q$ by
the seminorm Cauchy--Schwarz inequality.  Maximizing over $y$ gives
$r(x)\leq H-d^2/2+dR_Q\leq H+R_Q\sqrt{2H}$.
The lower inequality is Proposition~\ref{prop:theory-support}.
\end{proof}

\begin{proof}[Proof of Proposition~\ref{prop:exact-nominal}]
Let $\lambda$ be the one-bus mixture weight and $x$ the mean early load.
Complete-plan mixtures have $0\leq\lambda\leq1$,
$10\lambda\leq x\leq10$, and intrinsic cost $14-7\lambda$.
These conditions are sufficient: for $\lambda<1$ the two-bus component
buys $(x-10\lambda)/(1-\lambda)\in[0,10]$ early; at $\lambda=1$,
$x=10$.  For fixed $x$, cheapest intrinsic cost uses
$\lambda=x/10$.  Substituting $z=30-x$ into \eqref{eq:exact-supply}
gives hull objective
\[
 14-\frac{7x}{10}+F(x,30-x)
 =104-2.7x+0.2x^2
 =\frac{7591}{80}+\frac{(x-27/4)^2}{5}.
\]
Its minimizer is $x=27/4$, $\lambda=27/40$.  Direct minimization of
the one- and two-bus physical branches gives respectively 97 and 99.
At $(10,20)$ the gradient of $F$ is $(6,4)$; its private objective is
$7+6(10)+4(20)=147$, while the two-bus $x=0$ plan has
$14+4(30)=134$.  At hull load $(27/4,93/4)$ the gradient is
$(107/20,93/20)=(5.35,4.65)$.  Both supporting plans have private
objective 153.5, and $F^*(p)=58.6125$ on $\R_+^2$.
The hull dual value is therefore 94.8875.  The physical one-bus plan is
a price minimizer, so fleet LOC is zero; supplier LOC is $97-94.8875$.
\end{proof}

\begin{proof}[Proof of Proposition~\ref{prop:exact-robust}]
For reserve $0\leq r\leq5$, efficiency $\eta$, one-hour early limit
$P$, and one-hour single terminal connector $K$, write
\[
 T=30/\eta,\quad \ell=\max\{0,T-K\},\quad
 u=\min\{P,15/\eta\},\quad
 h=\max\{\ell,(10+r)/\eta\}.
\]
The two-bus physical early-energy interval is $[\ell,u]$, and the
one-bus interval is $[h,u]$ when nonempty.  The lower limit $\ell$
enforces terminal connector capacity, while $(10+r)/\eta$ ensures the
one bus can begin B above reserve.  For every point in these intervals,
charge only A's bus early; terminal demands total $T-x\leq K$ and can be
served serially in the one-hour terminal window.  Individual inventories
and acceptance limits hold, and all used buses finish at capacity.

Write $G_T(x)=F(x,T-x)$, $k=1/5$, and
$m=T/2-5a/2$ for early linear coefficient $a=4$.  Then
$G_T(x)=G_T(m)+k(x-m)^2$.  Suppose
$\ell<m<h<u$, $f>0$, and
$f<2k(h-\ell)(h-m)$.  The cheapest hull operating-cost boundary over
$[\ell,h]$ is $2f-f(x-\ell)/(h-\ell)$, obtained by mixing the
two-bus endpoint $\ell$ and one-bus endpoint $h$; over $[h,u]$ it is
$f$.  Put $d=h-\ell$.  The strict inequality places the hull minimum
$x^*=m+f/(2kd)$ inside $(m,h)$.  Physical branch minimization gives
\[
 D=G_T(m)+\min\{f+k(h-m)^2,2f\},\quad
 CH=G_T(m)+2f-\frac{f(m-\ell)}d-\frac{f^2}{4kd^2}.
\]
Consequently
\[
 \Delta=\min\left\{k(h-x^*)^2,
       \frac{f(m-\ell)}d+\frac{f^2}{4kd^2}\right\}>0.
\]
At $(r,\eta,P,K,f,a)=(1,19/20,12,30,7,4)$,
$(\ell,m,h,u,d)=(30/19,110/19,220/19,12,10)$.
The inequalities are strict: $u-h=8/19$ and
$2kd(h-m)=440/19>7$.  Rational substitution yields
\eqref{eq:exact-robust}.  All strict inequalities persist under small
changes in $(r,\eta,P)$ with $K=30$ fixed.
\end{proof}

\begin{proof}[Proof of Proposition~\ref{prop:exact-replication}]
For $m$ paired buses, total early energy $E$ ranges over
$[10m,10n]$, with intrinsic cost $14n-7m$ and terminal energy
$30n-E$.  Convexifying complete plans gives a triangular projected set;
its cheapest operating cost at each $E$ is $14n-0.7E$.  Substitution in
\eqref{eq:exact-scaled-supply} gives a quadratic minimized at
$E=27n/4$ with value $7591n/80$.

For fixed integer $m\geq n/2$, physical supply cost is minimized at
$E=10m$, and the total objective is
$104n-27m+20m^2/n=7591n/80+20(m-27n/40)^2/n$.
For $m\leq n/2$, the best early load is $E=5n$ when feasible and the
branch objective is $99n-7m$, decreasing toward $n/2$.  For odd $n$, moving from $m=\lfloor n/2\rfloor$ to
$m=\lceil n/2\rceil$ changes the minimized objective by $-7+5/n<0$.
For even $n$, the two formulas agree at $m=n/2$, and an integer nearest
to $27n/40$ lies in the upper regime.  Thus the nearest integer to
$27n/40$ is optimal,
with both integers at a half-integer tie and $|\delta|\leq1/2$.

At a physical optimum, the price is
$p=(4+2m/n,6-2m/n)$.  The one-bus private objective minus the best
two-bus private objective is $40\delta/n$; the latter buys zero early
energy when $m/n\geq1/2$.  Hence whole-operator regret is
\[
 r_n=\begin{cases}
 m\,40\delta/n,&\delta\geq0,\\
 (n-m)(-40\delta/n),&\delta<0.
 \end{cases}
\]
For $n=40k+1$, $m=27k+1$ and $\delta=13/40$, giving
\eqref{eq:exact-replication-regret}.  The same difference is a
reserved pair's individual regret when its current choice is worse;
$|40\delta/n|\leq20/n$.  Summing the disadvantaged pairs reproduces
the whole-operator expression because each pair has its own reserved
connector rights.
\end{proof}

\subsection{Averaged responses in the quadratic example}\label{app:averaged-responses}

Consider the exact recurrence in Section~\ref{sec:generation}, with
$\alpha_k=0.5/\sqrt{k}$ and projection onto $\R_+^2$. Put
$e_k=e(x_k)$, $c_k=c(x_k)$, $s_k=(5[p_{1,k}-4]_+,5p_{2,k})$,
$d_k=e_k-s_k$, and $A_K=\sum_{k=1}^K\alpha_k$.
The fleet loads satisfy $0\le e_{1,k}\le10$ and $20\le e_{2,k}\le30$.
For $k\ge7$, $5\alpha_k<1$. If $p_{1,k}\ge4$, its unprojected successor
is a convex combination of $p_{1,k}$ and $4+e_{1,k}/5\in[4,6]$.
If $p_{1,k}<4$, that successor is $p_{1,k}+\alpha_ke_{1,k}\in[0,6)$.
The second coordinate is a convex combination of $p_{2,k}$ and
$e_{2,k}/5\in[4,6]$. Thus projection is inactive from $k=7$ onward,
and prices, supply responses and $d_k$ remain bounded.

Summing $\alpha_kd_k=p_{k+1}-p_k$ from $k=7$ shows that
$\bar e_K-\bar s_K\to0$, where bars denote step-weighted averages.
The finite prefix has vanishing weight. Also,
\[
 \sum_{k=7}^K\alpha_kp_k^\top d_k
 =\frac12\left(\|p_{K+1}\|^2-\|p_7\|^2
                  -\sum_{k=7}^K\alpha_k^2\|d_k\|^2\right).
\]
Dividing by $A_K$ gives zero in the limit, since
$A_K\asymp\sqrt K$ and $\sum_{k=1}^K\alpha_k^2=O(\log K)$.
The projection inequality underlying \eqref{eq:dispatch-rate} gives the
stronger weighted statement
\[
 0\le\frac1{A_K}\sum_{k=1}^K\alpha_k[CH-Q(p_k)]
 \le\frac{\|p_1-p^*\|^2+M^2\sum_{k=1}^K\alpha_k^2}{2A_K}
 \longrightarrow0,
\]
with the attained price $p^*=(107/20,93/20)$ and a bound $M$ on
$\|d_k\|$. Exact oracle solutions satisfy
$c_k+F(s_k)=Q(p_k)-p_k^\top d_k$, so their step-weighted mean tends
to $CH$. Jensen's inequality, the vanishing load mismatch, and uniform
continuity of the quadratic on bounded loads now imply
\[
 CH\le\bar c_K+F(\bar e_K)
 \le\frac1{A_K}\sum_{k=1}^K\alpha_k[c_k+F(s_k)]+o(1)
 \longrightarrow CH.
\]
The averaged fleet point belongs to the compact complete-fleet hull.
Its unique optimizer therefore gives
$(\bar e_K,\bar c_K)\to((27/4,93/4),371/40)$.
Since $c_k$ equals 7 for one bus and 14 for two,
$\bar c_K=14-7\bar w_K$, where $\bar w_K$ is the step-weighted one-bus
frequency. Hence $\bar w_K\to27/40$. An eventually constant bus count
would instead give zero or one, so both counts occur infinitely often.

This proof concerns the stated quadratic recurrence and exact responses.
It does not apply a square-summable-step theorem to $1/\sqrt{k}$ steps,
or assert recovery of a physical fleet from its convexified average.
